\section{Introduction}

As Language Models (LMs) have scaled in complexity, the mechanistic interpretability community has developed a broad set of tools for studying the representations and algorithms underlying model behavior, including gradient-based attribution \citep{sundararajan2017axiomatic}, activation-based localization \citep{meng2022locating,meng2022mass}, and feature decomposition with Sparse Autoencoders (SAEs) \citep{bricken2023towards,cunningham2023sparse}. Among these paradigms, causal abstraction provides a particularly rigorous framework: given a task, a high-level causal model and an alignment between high-level variables and internal neural representations, it asks whether the model and the hypothesis agree under counterfactual interventions \citep{geiger2021causal, geiger2025causal}. This framework has been shown to be suitable for a wide variety of tasks \citep{wu2024pyvene, arora-etal-2024-causalgym, huanginternal, boguraev-etal-2025-causal}.

In practice, causal abstraction is usually evaluated through \emph{interchange interventions}, summarized by a single scalar metric, \emph{interchange intervention accuracy} (IIA) \citep{geiger2021causal, geiger2025causal}. The IIA score provides a simple metric to measure how well hypotheses align with neural models, but it says little about \emph{where} a hypothesis is faithful, and intermediate scores are especially hard to interpret \citep{makelov2023subspace,wu2024reply,meloux2025everything}. Moreover, although methods such as DAS make it easier to find candidate alignments, the overall workflow remains largely evaluative rather than constructive: it can tell us that a hypothesis is imperfect, but not how to improve it.

We address this gap by shifting the unit of analysis from a global score to the structure of the input space. Given a low-level model \(\mathcal L\), a high-level hypothesis \(\mathcal H\), and an alignment \(\Pi\), we partition the input space into well-interpreted \emph{target buckets} with high IIA and a complementary bucket that captures the remaining failure modes. The key question is no longer just whether an abstraction works on average, but \emph{for which inputs} it is actually faithful. This yields a more informative diagnosis and turns abstraction failures into evidence for refinement. Similar ideas of exploiting structure in observational or interventional partitions have appeared in causal feature learning \citep{chalupka2014visual,chalupka2017causal}.

Our main contribution is a practical four-step pipeline for diagnosing and improving causal abstractions: specify a reliable task and input space, obtain a candidate alignment, partition the input space by pairwise interchangeability to identify nearly interchange-consistent subsets, and train a classifier to generalize this diagnosis beyond the analyzed sample. Across fine-tuned and pretrained models of different sizes, and across alignment methods including full-vector patching, DAS, and MDAS, we show that this bucketing procedure is broadly useful for diagnosing causal abstractions. Empirically, we validate it on logic, entity binding \citep{gur2025mixing}, and factual recall \citep{geva2021transformer,hernandez2023linearity,geva2023dissecting,huang2024ravel}. In the toy logic task, recursively applying the recipe supports iterative refinement of the high-level hypothesis itself, showing how causal abstraction can move from post hoc evaluation toward constructive hypothesis discovery.

\begin{figure}[t]
    \centering
    \includegraphics[width=\linewidth]{figs/teaser_new.png}
    \caption{\textbf{Four-step interpretation diagnosis pipeline.} Given a causal abstraction for a task-performing model, we focus on task-correct inputs, identify an alignment, partition the input space by pairwise interchangeability, and train a classifier to generalize the diagnosis.}
    \label{fig:teaser}
\end{figure}

The rest of the paper is organized as follows. Section~2 reviews the causal abstraction framework and the alignment and feature-learning tools used in this work. Section~3 formalizes interchange-consistent subsets and introduces our diagnosis pipeline. Section~4 evaluates the method on three settings of increasing complexity, and Section~5 concludes with limitations and future directions.